\originalsection{18.04第二次模拟考试}
\sourceinfo{Jeremy Orloff, Spring 2018. 题面 mit18\_04s18\_pex2.pdf, PDF pp.1--2; 官方解答 pex2-qa, pp.1--11. MIT OCW, CC BY-NC-SA 4.0. 12道大题.}
\begin{exercise}\label{ass:pex2-1}
\textbf{原题 Pex2.1。} 调和函数均指实值函数. (a) 证明 $u=x^3-3xy^2+3x^2-3y^2$ 调和并求调和共轭. (b) 求单位圆盘内满足 $u(1/2)=2$ 且处处 $u\ge2$ 的所有调和函数. (c) 单位圆周四象限的温度依次固定为1,2,3,4, 求圆心温度. (d) 共轭调和函数 $u,v$ 的乘积 $uv$ 是否调和? (e) 证明 $u$ 调和则 $u_x$ 调和. (f) 证明 $u$ 与 $u^2$ 均调和时 $u$ 常数.
\end{exercise}
\begin{solution}
据官方解答翻译并补全正则性及边界解释. (a) $u_{xx}=6x+6,u_{yy}=-6x-6$, 和为0. $u=\Re(z^3+3z^2)$, 故共轭为 $v=3x^2y-y^3+6xy+C$, $C\in\R$; 也可积分 $v_y=u_x,v_x=-u_y$ 得到.

(b) $u$ 在内部点取最小值, 强最小值原理给 $u\equiv2$. (c) 以 Poisson 积分取具有这些分段边界值的有界调和温度; 象限端点单点值不影响积分. 圆心是边界平均, 即 $(1+2+3+4)/4=5/2$. (d) $f=u+\ii v$ 解析, $\Im(f^2)=2uv$, 所以 $uv$ 调和. (e) 调和函数内部光滑, 故 $\Delta u_x=\partial_x\Delta u=0$; 在局部选调和共轭, $u_x=\Re f'$ 也给出同样结论, 不必假定全区域单连通. (f) $0=\Delta(u^2)=2(u_x^2+u_y^2)+2u\Delta u$, 故梯度为0, 在每个连通分支上 $u$ 常数; 按课程“区域”约定即全域常数.
\end{solution}
\begin{exercise}\label{ass:pex2-2}
\textbf{原题 Pex2.2。} 求 $f(z)=1/((z-1)(z-3))$ 在以0为中心、使 $f$ 解析的三个环域内的 Laurent 级数.
\end{exercise}
\begin{solution}
据官方解答翻译. 部分分式为 $f=-1/(2(z-1))+1/(2(z-3))$. 按内外几何级数分别展开, 得
\[
\begin{aligned}
|z|<1:\quad&f=\frac12\sum_{n=0}^{\infty}(1-3^{-n-1})z^n,\\
1<|z|<3:\quad&f=-\frac12\sum_{n=1}^{\infty}z^{-n}-\frac16\sum_{n=0}^{\infty}(z/3)^n,\\
|z|>3:\quad&f=\frac12\sum_{n=1}^{\infty}(3^{n-1}-1)z^{-n}.
\end{aligned}
\]
例如 $1/(z-a)=-a^{-1}\sum(z/a)^n$ 在内侧, $z^{-1}\sum(a/z)^n$ 在外侧, 各自收敛条件给上述三个区域. 0本来是普通点, 第一区域含0.
\end{solution}
\begin{exercise}\label{ass:pex2-3}
\textbf{原题 Pex2.3。} 求0附近 Laurent 级数的前几项: (a) $z^2\cos(1/(3z))$, $0<|z|$; (b) $1/(\ee^z-1)$, $0<|z|<R$, 并确定 $R$.
\end{exercise}
\begin{solution}
据官方解答翻译; (b)的最大收敛半径是编者补充（AI）, 原答案未给. (a) 代入余弦级数得 $z^2-1/18+1/(1944z^2)-1/(524880z^4)+\cdots$, 通项 $(-1)^n z^{2-2n}/(3^{2n}(2n)!)$, 在任意 $z\ne0$ 收敛. (b) 将 $\ee^z-1=z(1+z/2+z^2/6+\cdots)$, 倒数系数逐次匹配得 $1/z-1/2+z/12-z^3/720+\cdots$. 最近的其他奇点为 $\pm2\pi\ii$, 所以最大环域为 $0<|z|<2\pi$, 即 $R=2\pi$.
\end{solution}
\begin{exercise}\label{ass:pex2-4}
\textbf{原题 Pex2.4。} 求 $\sum_{n=-\infty}^{\infty}z^n/2^{|n|}$ 的收敛环域.
\end{exercise}
\begin{solution}据官方解答翻译. 非负次部分为 $\sum_{n\ge0}(z/2)^n$, 收敛要求 $|z|<2$; 负次部分为 $\sum_{k\ge1}(1/(2z))^k$, 要求 $|z|>1/2$. 合起来 $1/2<|z|<2$. 两边界至少一部分的通项模不趋0, 均不收敛.\end{solution}
\begin{exercise}\label{ass:pex2-5}
\textbf{原题 Pex2.5。} 找出、分类各孤立奇点并求留数: (a) $(z^3+1)/(z^2(z+1))$; (b) $1/(\ee^z-1)$; (c) $\cos(1-1/z)$.
\end{exercise}
\begin{solution}
据官方解答翻译并补全可去性的理由. (a) 约分为 $1-z^{-1}+z^{-2}$, 故0是二阶极点, 留数 $-1$; 原式 $-1$ 处可去, 留数0, 延拓值3. (b) $2k\pi\ii$ ($k\in\Z$) 为全部奇点, 分母在每点一阶零, 留数 $1/\ee^{2k\pi\ii}=1$. (c) 唯一有限奇点为0, 因为 $\cos(1-1/z)=\cos1\cos(1/z)+\sin1\sin(1/z)$, 无穷多个负幂系数非零, 故本性奇点. $z^{-1}$ 项只来自 $\sin1\sin(1/z)$, 留数为 $\sin1$.
\end{solution}
\begin{exercise}\label{ass:pex2-6}
\textbf{原题 Pex2.6。} 构造函数: (a) $1+\ii$ 为二阶极点, 0、1为本性奇点; (b) 0为可去奇点, 1为六阶极点, $\ii$ 为本性奇点.
\end{exercise}
\begin{solution}
(a) 据官方解答, 取 $\ee^{1/z}+\ee^{1/(z-1)}+(z-1-\ii)^{-2}$. 每个指定点仅一个加项奇异, 其余两项在该点解析, 所以不会抵消指定主部.

(b) 编者补充（AI）的一种较简单构造为 $\sin z/z+(z-1)^{-6}+\ee^{1/(z-\ii)}$, 初始定义时删去0、1、$\ii$. 第一项0处可去, 后两项0处解析; 1处只有第二项奇异, 为六阶极点; $\ii$ 处只有第三项有无穷主部, 为本性奇点. 原官方解答另取 $(z^2+8z)/\sin z$ 作第一项, 也能满足指定点条件, 但引入其他极点.
\end{solution}
\begin{exercise}\label{ass:pex2-7}
\textbf{原题 Pex2.7。} 判断真伪; 真则证明、假则给反例. (a) $f,g$ 在 $z_0$ 均有极点, 则 $f+g$ 在该点有极点. (b) 若两者留数均非零, 则 $fg$ 有极点且留数非零. (c) $f$ 在0本性奇异而 $g$ 在0有有限阶极点, 则 $f+g$ 在0本性奇异. (d) $f$ 在0有 $m$ 阶极点, 则 $f(z^2)$ 在0有 $2m$ 阶极点.
\end{exercise}
\begin{solution}
据官方解答翻译并校正(d)的记号与“零点”误写. (a) 假, 取 $z_0=0,f=1/z,g=-1/z$, 和恒0. (b) 假, $f=g=1/z$ 各留数1, 乘积 $1/z^2$ 留数0. 乘积确有极点, 错在其留数断言. (c) 真, $g$ 只有有限多个负幂项, 不能消除 $f$ 无穷多个非零负幂系数. (d) 真, $f(z)=z^{-m}h(z)$, $h$ 解析且 $h(0)\ne0$, 所以 $f(z^2)=z^{-2m}h(z^2)$ 为 $2m$ 阶极点. 官方把左侧误写为 $f(z)^2$ 且最后写成 $2m$ 阶零点, 均已明确修正.
\end{solution}
\begin{exercise}\label{ass:pex2-8}
\textbf{原题 Pex2.8(a)。} 在 $|z|>1$ 求 $1/\ee^{1-z}$ 的 Laurent 级数. 原卷第8题只列(a).
\end{exercise}
\begin{solution}据官方解答翻译. $\ee^{z-1}=\ee^{-1}\sum_{n=0}^{\infty}z^n/n!$, 此 Taylor 级数就是所求 Laurent 级数, 无负幂项. 它实际全平面收敛, 当然也适用于指定外部区域.\end{solution}
\begin{exercise}\label{ass:pex2-9}
\textbf{原题 Pex2.9。} 在 $|z|<2\pi$ 令 $h(z)=1/\sin z-1/z+2z/(z^2-\pi^2)$. (a) 证明所有表观奇点均可去. (b) 求0附近 Taylor 级数的前4项.
\end{exercise}
\begin{solution}
据官方解答翻译, (b)补足其未算出的项（AI）. (a) 仅需检查 $0,\pm\pi$. $1/\sin z$ 在0留数1, 被 $-1/z$ 消去; 在 $\pm\pi$ 留数均 $-1$. 因 $2z/(z^2-\pi^2)=1/(z-\pi)+1/(z+\pi)$, 另两个主部也消去. 各点原本只是简单极点, 故无负幂后都可去.

(b) 由 $\csc z=z^{-1}+z/6+7z^3/360+31z^5/15120+127z^7/604800+\cdots$ 及 $2z/(z^2-\pi^2)=-\sum_{n\ge0}2z^{2n+1}/\pi^{2n+2}$, 前4个非零项为
\[
h(z)=\left(\frac16-\frac2{\pi^2}\right)z+
\left(\frac7{360}-\frac2{\pi^4}\right)z^3+
\left(\frac{31}{15120}-\frac2{\pi^6}\right)z^5+
\left(\frac{127}{604800}-\frac2{\pi^8}\right)z^7+O(z^9).
\]
第一展开可用 $\sin z$ 的幂级数乘法逐次验系数. 官方只算到 $z^3$, 并提常数项为0, 未达到四个非零项.
\end{solution}
\begin{exercise}\label{ass:pex2-10}
\textbf{原题 Pex2.10。} 求无穷远处留数: (a) $\ee^z$; (b) $(z-1)/(z+1)$.
\end{exercise}
\begin{solution}
据官方解答翻译. 留数在孤立本性奇点也有定义. (a) $\Res_\infty\ee^z=-\Res_{w=0}(\ee^{1/w}/w^2)=0$, 因后者各幂次数为 $-2,-3,\ldots$, 无 $w^{-1}$. (b) $(z-1)/(z+1)=1-2/z+2/z^2-\cdots$ 在 $|z|>1$ 成立, 所以无穷远留数是 $z^{-1}$ 系数的相反数2. 也可用有限唯一极点 $-1$ 的留数 $-2$ 验证.
\end{solution}
\begin{exercise}\label{ass:pex2-11}
\textbf{原题 Pex2.11。} 复势 $\Phi(z)=z+\log(z-\ii)+\log(z+\ii)$. 依次 (i) 辨认源、汇等流动成分; (ii) 求停滞点; (iii) 画每个源附近流动; (iv) 画远处流动; (v) 合成流线图.
\end{exercise}
\begin{solution}
据官方解答翻译, 流线数值重绘与局部结构解释为编者补充（AI）. (i) $z$ 表示向右的单位匀流; 两个系数为1的对数分别在 $\pm\ii$ 表示等强点源, 没有点汇. 对数需局部选择分支, 但速度 $\Phi'=1+2z/(z^2+1)=(z+1)^2/(z^2+1)$ 单值.

(ii) 唯一停滞位置是 $z=-1$, 导数在此二阶零, 不是两个不同停滞点. (iii) 在源 $a=\pm\ii$ 附近 $\Phi'=1/(z-a)+O(1)$, 实速度为 $\overline{\Phi'}\sim1/(\bar z-\bar a)$, 沿径向向外. (iv) 无穷远 $\Phi'=1+2/z+O(z^{-3})$, 故流动趋于向右匀流. (v) 流线满足 $\Im\Phi$ 局部常数, 实轴在源外也是流线. 图中积分的是实速度方向, 在源与停滞点附近截断; 箭头用于指示方向, 该采样图不充当理论证明. $-1$ 附近 $\Phi-\Phi(-1)\sim(z+1)^3/6$, 给出六条等流函数射线的局部结构.
\begin{center}\begin{tikzpicture}[xscale=1.35,yscale=1.25,>=Stealth]
\clip (-4.3,-3.3) rectangle (4.7,3.4);
\draw[main!65,line width=.35pt] plot[smooth] coordinates {(0.1200,1.0000) (0.1450,1.0016) (0.2444,1.0120) (0.3432,1.0276) (0.4413,1.0468) (0.5389,1.0685) (0.6361,1.0919) (0.7331,1.1163) (0.8299,1.1413) (0.9267,1.1665) (1.0235,1.1917) (1.1203,1.2168) (1.2172,1.2416) (1.3141,1.2659) (1.4112,1.2898) (1.5085,1.3132) (1.6058,1.3360) (1.7033,1.3583) (1.8010,1.3799) (1.8987,1.4010) (1.9966,1.4215) (2.0946,1.4414) (2.1927,1.4608) (2.2909,1.4796) (2.3892,1.4978) (2.4876,1.5155) (2.5862,1.5327) (2.6848,1.5494) (2.7834,1.5657) (2.8822,1.5814) (2.9810,1.5967) (3.0799,1.6116) (3.1788,1.6260) (3.2779,1.6400) (3.3769,1.6536) (3.4760,1.6669) (3.5752,1.6798) (3.6744,1.6923) (3.7737,1.7045) (3.8730,1.7163) (3.9723,1.7279)};
\draw[main!65,line width=.35pt] plot[smooth] coordinates {(0.1039,1.0600) (0.1257,1.0722) (0.2134,1.1203) (0.3019,1.1668) (0.3912,1.2119) (0.4812,1.2554) (0.5719,1.2975) (0.6633,1.3382) (0.7553,1.3774) (0.8478,1.4153) (0.9409,1.4519) (1.0344,1.4872) (1.1285,1.5213) (1.2229,1.5542) (1.3177,1.5860) (1.4129,1.6167) (1.5084,1.6464) (1.6042,1.6751) (1.7002,1.7028) (1.7966,1.7296) (1.8932,1.7555) (1.9900,1.7806) (2.0870,1.8048) (2.1842,1.8283) (2.2816,1.8511) (2.3791,1.8731) (2.4768,1.8944) (2.5746,1.9151) (2.6726,1.9352) (2.7707,1.9547) (2.8689,1.9735) (2.9672,1.9919) (3.0656,2.0096) (3.1641,2.0269) (3.2627,2.0437) (3.3613,2.0600) (3.4601,2.0759) (3.5589,2.0913) (3.6577,2.1063) (3.7567,2.1208) (3.8557,2.1350) (3.9547,2.1488)};
\draw[main!65,line width=.35pt] plot[smooth] coordinates {(0.0600,1.1039) (0.0742,1.1245) (0.1359,1.2031) (0.2039,1.2765) (0.2766,1.3450) (0.3532,1.4093) (0.4329,1.4698) (0.5150,1.5269) (0.5991,1.5808) (0.6850,1.6320) (0.7724,1.6807) (0.8610,1.7271) (0.9506,1.7714) (1.0412,1.8137) (1.1327,1.8542) (1.2248,1.8931) (1.3176,1.9304) (1.4110,1.9662) (1.5048,2.0006) (1.5992,2.0338) (1.6939,2.0658) (1.7890,2.0966) (1.8845,2.1264) (1.9803,2.1551) (2.0764,2.1828) (2.1727,2.2097) (2.2693,2.2356) (2.3661,2.2607) (2.4631,2.2851) (2.5602,2.3087) (2.6576,2.3315) (2.7551,2.3537) (2.8528,2.3752) (2.9506,2.3960) (3.0485,2.4163) (3.1465,2.4359) (3.2447,2.4551) (3.3430,2.4736) (3.4413,2.4917) (3.5398,2.5093) (3.6383,2.5264) (3.7369,2.5430) (3.8356,2.5592) (3.9343,2.5750)};
\draw[main!65,line width=.35pt] plot[smooth] coordinates {(0.0000,1.1200) (0.0034,1.1448) (0.0267,1.2419) (0.0635,1.3348) (0.1111,1.4227) (0.1670,1.5056) (0.2296,1.5835) (0.2976,1.6569) (0.3698,1.7260) (0.4456,1.7912) (0.5243,1.8529) (0.6054,1.9114) (0.6886,1.9669) (0.7735,2.0197) (0.8599,2.0700) (0.9476,2.1181) (1.0364,2.1640) (1.1262,2.2080) (1.2169,2.2502) (1.3083,2.2906) (1.4005,2.3295) (1.4932,2.3669) (1.5865,2.4030) (1.6803,2.4377) (1.7745,2.4712) (1.8691,2.5035) (1.9641,2.5347) (2.0595,2.5649) (2.1551,2.5941) (2.2510,2.6223) (2.3472,2.6497) (2.4436,2.6762) (2.5403,2.7019) (2.6371,2.7268) (2.7341,2.7510) (2.8314,2.7745) (2.9287,2.7973) (3.0262,2.8195) (3.1239,2.8410) (3.2217,2.8620) (3.3196,2.8824) (3.4176,2.9022) (3.5157,2.9215) (3.6139,2.9403) (3.7122,2.9586) (3.8106,2.9764) (3.9091,2.9938)};
\draw[main!65,line width=.35pt] plot[smooth] coordinates {(-0.0600,1.1039) (-0.0689,1.1273) (-0.0923,1.2243) (-0.0968,1.3241) (-0.0848,1.4233) (-0.0589,1.5198) (-0.0216,1.6125) (0.0250,1.7009) (0.0792,1.7850) (0.1395,1.8647) (0.2050,1.9403) (0.2746,2.0120) (0.3479,2.0801) (0.4241,2.1448) (0.5028,2.2065) (0.5837,2.2652) (0.6665,2.3213) (0.7509,2.3749) (0.8368,2.4262) (0.9239,2.4753) (1.0120,2.5225) (1.1012,2.5678) (1.1912,2.6113) (1.2820,2.6532) (1.3735,2.6935) (1.4657,2.7324) (1.5584,2.7699) (1.6516,2.8062) (1.7452,2.8412) (1.8393,2.8750) (1.9338,2.9077) (2.0287,2.9395) (2.1238,2.9702) (2.2193,2.9999)};
\draw[main!65,line width=.35pt] plot[smooth] coordinates {(-0.1039,1.0600) (-0.1236,1.0754) (-0.1933,1.1469) (-0.2462,1.2315) (-0.2808,1.3252) (-0.2973,1.4237) (-0.2975,1.5236) (-0.2835,1.6225) (-0.2577,1.7191) (-0.2218,1.8124) (-0.1776,1.9020) (-0.1265,1.9880) (-0.0695,2.0701) (-0.0076,2.1486) (0.0586,2.2236) (0.1283,2.2953) (0.2011,2.3638) (0.2767,2.4293) (0.3545,2.4921) (0.4344,2.5522) (0.5161,2.6098) (0.5994,2.6652) (0.6841,2.7184) (0.7700,2.7695) (0.8570,2.8187) (0.9451,2.8661) (1.0340,2.9118) (1.1238,2.9559) (1.2143,2.9985)};
\draw[main!65,line width=.35pt] plot[smooth] coordinates {(-0.1200,1.0000) (-0.1449,1.0021) (-0.2433,1.0193) (-0.3373,1.0530) (-0.4227,1.1048) (-0.4948,1.1737) (-0.5506,1.2565) (-0.5890,1.3487) (-0.6108,1.4461) (-0.6175,1.5458) (-0.6112,1.6456) (-0.5938,1.7440) (-0.5668,1.8403) (-0.5317,1.9339) (-0.4895,2.0245) (-0.4413,2.1121) (-0.3878,2.1966) (-0.3297,2.2780) (-0.2676,2.3563) (-0.2020,2.4318) (-0.1333,2.5044) (-0.0617,2.5743) (0.0122,2.6415) (0.0884,2.7063) (0.1665,2.7688) (0.2464,2.8289) (0.3279,2.8869) (0.4107,2.9429) (0.4949,2.9969)};
\draw[main!65,line width=.35pt] plot[smooth] coordinates {(-0.1039,0.9400) (-0.1255,0.9274) (-0.2116,0.8765) (-0.2977,0.8256) (-0.3843,0.7757) (-0.4723,0.7283) (-0.5629,0.6859) (-0.6573,0.6532) (-0.7559,0.6376) (-0.8550,0.6480) (-0.9457,0.6891) (-1.0200,0.7555) (-1.0762,0.8380) (-1.1164,0.9294) (-1.1434,1.0257) (-1.1594,1.1243) (-1.1661,1.2241) (-1.1645,1.3240) (-1.1555,1.4236) (-1.1399,1.5224) (-1.1182,1.6200) (-1.0909,1.7161) (-1.0583,1.8107) (-1.0209,1.9034) (-0.9790,1.9942) (-0.9329,2.0830) (-0.8830,2.1696) (-0.8294,2.2540) (-0.7725,2.3362) (-0.7125,2.4162) (-0.6496,2.4939) (-0.5840,2.5695) (-0.5161,2.6428) (-0.4458,2.7140) (-0.3735,2.7830) (-0.2992,2.8500) (-0.2232,2.9149) (-0.1455,2.9779)};
\draw[main!65,line width=.35pt] plot[smooth] coordinates {(-0.0600,0.8961) (-0.0703,0.8733) (-0.1020,0.7786) (-0.1147,0.6796) (-0.1034,0.5805) (-0.0661,0.4881) (-0.0051,0.4092) (0.0729,0.3470) (0.1614,0.3007) (0.2557,0.2674) (0.3528,0.2439) (0.4515,0.2276) (0.5508,0.2165) (0.6506,0.2092) (0.7505,0.2048) (0.8504,0.2023) (0.9504,0.2014) (1.0504,0.2016) (1.1504,0.2026) (1.2504,0.2043) (1.3504,0.2063) (1.4504,0.2087) (1.5503,0.2114) (1.6503,0.2142) (1.7503,0.2171) (1.8502,0.2201) (1.9502,0.2231) (2.0501,0.2261) (2.1501,0.2291) (2.2500,0.2320) (2.3500,0.2350) (2.4499,0.2379) (2.5499,0.2407) (2.6499,0.2435) (2.7498,0.2462) (2.8498,0.2488) (2.9498,0.2514) (3.0497,0.2539) (3.1497,0.2564) (3.2497,0.2588) (3.3496,0.2611) (3.4496,0.2634) (3.5496,0.2656) (3.6496,0.2678) (3.7495,0.2699) (3.8495,0.2719) (3.9495,0.2739)};
\draw[main!65,line width=.35pt] plot[smooth] coordinates {(-0.0000,0.8800) (0.0039,0.8553) (0.0336,0.7601) (0.0852,0.6747) (0.1551,0.6035) (0.2381,0.5479) (0.3291,0.5067) (0.4247,0.4775) (0.5226,0.4575) (0.6218,0.4447) (0.7215,0.4371) (0.8214,0.4334) (0.9214,0.4327) (1.0214,0.4341) (1.1213,0.4371) (1.2213,0.4413) (1.3211,0.4463) (1.4210,0.4519) (1.5208,0.4580) (1.6206,0.4644) (1.7204,0.4709) (1.8201,0.4776) (1.9199,0.4843) (2.0197,0.4910) (2.1195,0.4976) (2.2193,0.5042) (2.3191,0.5106) (2.4189,0.5170) (2.5187,0.5232) (2.6185,0.5293) (2.7183,0.5352) (2.8181,0.5410) (2.9180,0.5467) (3.0178,0.5522) (3.1177,0.5576) (3.2175,0.5629) (3.3174,0.5680) (3.4173,0.5729) (3.5172,0.5778) (3.6170,0.5825) (3.7169,0.5871) (3.8168,0.5915) (3.9168,0.5959) (4.0167,0.6001)};
\draw[main!65,line width=.35pt] plot[smooth] coordinates {(0.0600,0.8961) (0.0759,0.8767) (0.1489,0.8087) (0.2338,0.7561) (0.3261,0.7179) (0.4226,0.6919) (0.5212,0.6756) (0.6208,0.6666) (0.7208,0.6632) (0.8207,0.6640) (0.9207,0.6679) (1.0205,0.6740) (1.1202,0.6817) (1.2198,0.6906) (1.3193,0.7004) (1.4188,0.7107) (1.5182,0.7214) (1.6176,0.7324) (1.7170,0.7434) (1.8164,0.7544) (1.9158,0.7654) (2.0152,0.7762) (2.1146,0.7869) (2.2141,0.7974) (2.3135,0.8077) (2.4130,0.8178) (2.5125,0.8276) (2.6121,0.8373) (2.7116,0.8467) (2.8112,0.8558) (2.9108,0.8647) (3.0104,0.8734) (3.1101,0.8819) (3.2097,0.8901) (3.3094,0.8981) (3.4091,0.9059) (3.5088,0.9135) (3.6085,0.9208) (3.7083,0.9280) (3.8080,0.9350) (3.9078,0.9418) (4.0076,0.9484)};
\draw[main!65,line width=.35pt] plot[smooth] coordinates {(0.1039,0.9400) (0.1270,0.9303) (0.2223,0.9004) (0.3207,0.8826) (0.4203,0.8742) (0.5203,0.8731) (0.6202,0.8773) (0.7198,0.8854) (0.8192,0.8964) (0.9184,0.9095) (1.0173,0.9239) (1.1161,0.9394) (1.2148,0.9554) (1.3134,0.9719) (1.4121,0.9885) (1.5107,1.0051) (1.6093,1.0217) (1.7079,1.0381) (1.8066,1.0542) (1.9054,1.0701) (2.0041,1.0856) (2.1030,1.1009) (2.2019,1.1157) (2.3008,1.1302) (2.3998,1.1444) (2.4988,1.1581) (2.5979,1.1715) (2.6971,1.1846) (2.7963,1.1972) (2.8955,1.2096) (2.9948,1.2216) (3.0941,1.2332) (3.1935,1.2446) (3.2929,1.2556) (3.3923,1.2663) (3.4917,1.2768) (3.5912,1.2869) (3.6907,1.2968) (3.7903,1.3064) (3.8898,1.3157) (3.9894,1.3248)};
\draw[main!65,line width=.35pt] plot[smooth] coordinates {(0.1200,-1.0000) (0.1450,-1.0016) (0.2444,-1.0120) (0.3432,-1.0276) (0.4413,-1.0468) (0.5389,-1.0685) (0.6361,-1.0919) (0.7331,-1.1163) (0.8299,-1.1413) (0.9267,-1.1665) (1.0235,-1.1917) (1.1203,-1.2168) (1.2172,-1.2416) (1.3141,-1.2659) (1.4112,-1.2898) (1.5085,-1.3132) (1.6058,-1.3360) (1.7033,-1.3583) (1.8010,-1.3799) (1.8987,-1.4010) (1.9966,-1.4215) (2.0946,-1.4414) (2.1927,-1.4608) (2.2909,-1.4796) (2.3892,-1.4978) (2.4876,-1.5155) (2.5862,-1.5327) (2.6848,-1.5494) (2.7834,-1.5657) (2.8822,-1.5814) (2.9810,-1.5967) (3.0799,-1.6116) (3.1788,-1.6260) (3.2779,-1.6400) (3.3769,-1.6536) (3.4760,-1.6669) (3.5752,-1.6798) (3.6744,-1.6923) (3.7737,-1.7045) (3.8730,-1.7163) (3.9723,-1.7279)};
\draw[main!65,line width=.35pt] plot[smooth] coordinates {(0.1039,-0.9400) (0.1270,-0.9303) (0.2223,-0.9004) (0.3207,-0.8826) (0.4203,-0.8742) (0.5203,-0.8731) (0.6202,-0.8773) (0.7198,-0.8854) (0.8192,-0.8964) (0.9184,-0.9095) (1.0173,-0.9239) (1.1161,-0.9394) (1.2148,-0.9554) (1.3134,-0.9719) (1.4121,-0.9885) (1.5107,-1.0051) (1.6093,-1.0217) (1.7079,-1.0381) (1.8066,-1.0542) (1.9054,-1.0701) (2.0041,-1.0856) (2.1030,-1.1009) (2.2019,-1.1157) (2.3008,-1.1302) (2.3998,-1.1444) (2.4988,-1.1581) (2.5979,-1.1715) (2.6971,-1.1846) (2.7963,-1.1972) (2.8955,-1.2096) (2.9948,-1.2216) (3.0941,-1.2332) (3.1935,-1.2446) (3.2929,-1.2556) (3.3923,-1.2663) (3.4917,-1.2768) (3.5912,-1.2869) (3.6907,-1.2968) (3.7903,-1.3064) (3.8898,-1.3157) (3.9894,-1.3248)};
\draw[main!65,line width=.35pt] plot[smooth] coordinates {(0.0600,-0.8961) (0.0759,-0.8767) (0.1489,-0.8087) (0.2338,-0.7561) (0.3261,-0.7179) (0.4226,-0.6919) (0.5212,-0.6756) (0.6208,-0.6666) (0.7208,-0.6632) (0.8207,-0.6640) (0.9207,-0.6679) (1.0205,-0.6740) (1.1202,-0.6817) (1.2198,-0.6906) (1.3193,-0.7004) (1.4188,-0.7107) (1.5182,-0.7214) (1.6176,-0.7324) (1.7170,-0.7434) (1.8164,-0.7544) (1.9158,-0.7654) (2.0152,-0.7762) (2.1146,-0.7869) (2.2141,-0.7974) (2.3135,-0.8077) (2.4130,-0.8178) (2.5125,-0.8276) (2.6121,-0.8373) (2.7116,-0.8467) (2.8112,-0.8558) (2.9108,-0.8647) (3.0104,-0.8734) (3.1101,-0.8819) (3.2097,-0.8901) (3.3094,-0.8981) (3.4091,-0.9059) (3.5088,-0.9135) (3.6085,-0.9208) (3.7083,-0.9280) (3.8080,-0.9350) (3.9078,-0.9418) (4.0076,-0.9484)};
\draw[main!65,line width=.35pt] plot[smooth] coordinates {(0.0000,-0.8800) (0.0039,-0.8553) (0.0336,-0.7601) (0.0852,-0.6747) (0.1551,-0.6035) (0.2381,-0.5479) (0.3291,-0.5067) (0.4247,-0.4775) (0.5226,-0.4575) (0.6218,-0.4447) (0.7215,-0.4371) (0.8214,-0.4334) (0.9214,-0.4327) (1.0214,-0.4341) (1.1213,-0.4371) (1.2213,-0.4413) (1.3211,-0.4463) (1.4210,-0.4519) (1.5208,-0.4580) (1.6206,-0.4644) (1.7204,-0.4709) (1.8201,-0.4776) (1.9199,-0.4843) (2.0197,-0.4910) (2.1195,-0.4976) (2.2193,-0.5042) (2.3191,-0.5106) (2.4189,-0.5170) (2.5187,-0.5232) (2.6185,-0.5293) (2.7183,-0.5352) (2.8181,-0.5410) (2.9180,-0.5467) (3.0178,-0.5522) (3.1177,-0.5576) (3.2175,-0.5629) (3.3174,-0.5680) (3.4173,-0.5729) (3.5172,-0.5778) (3.6170,-0.5825) (3.7169,-0.5871) (3.8168,-0.5915) (3.9168,-0.5959) (4.0167,-0.6001)};
\draw[main!65,line width=.35pt] plot[smooth] coordinates {(-0.0600,-0.8961) (-0.0703,-0.8733) (-0.1020,-0.7786) (-0.1147,-0.6796) (-0.1034,-0.5805) (-0.0661,-0.4881) (-0.0051,-0.4092) (0.0729,-0.3470) (0.1614,-0.3007) (0.2557,-0.2674) (0.3528,-0.2439) (0.4515,-0.2276) (0.5508,-0.2165) (0.6506,-0.2092) (0.7505,-0.2048) (0.8504,-0.2023) (0.9504,-0.2014) (1.0504,-0.2016) (1.1504,-0.2026) (1.2504,-0.2043) (1.3504,-0.2063) (1.4504,-0.2087) (1.5503,-0.2114) (1.6503,-0.2142) (1.7503,-0.2171) (1.8502,-0.2201) (1.9502,-0.2231) (2.0501,-0.2261) (2.1501,-0.2291) (2.2500,-0.2320) (2.3500,-0.2350) (2.4499,-0.2379) (2.5499,-0.2407) (2.6499,-0.2435) (2.7498,-0.2462) (2.8498,-0.2488) (2.9498,-0.2514) (3.0497,-0.2539) (3.1497,-0.2564) (3.2497,-0.2588) (3.3496,-0.2611) (3.4496,-0.2634) (3.5496,-0.2656) (3.6496,-0.2678) (3.7495,-0.2699) (3.8495,-0.2719) (3.9495,-0.2739)};
\draw[main!65,line width=.35pt] plot[smooth] coordinates {(-0.1039,-0.9400) (-0.1255,-0.9274) (-0.2116,-0.8765) (-0.2977,-0.8256) (-0.3843,-0.7757) (-0.4723,-0.7283) (-0.5629,-0.6859) (-0.6573,-0.6532) (-0.7559,-0.6376) (-0.8550,-0.6480) (-0.9457,-0.6891) (-1.0200,-0.7555) (-1.0762,-0.8380) (-1.1164,-0.9294) (-1.1434,-1.0257) (-1.1594,-1.1243) (-1.1661,-1.2241) (-1.1645,-1.3240) (-1.1555,-1.4236) (-1.1399,-1.5224) (-1.1182,-1.6200) (-1.0909,-1.7161) (-1.0583,-1.8107) (-1.0209,-1.9034) (-0.9790,-1.9942) (-0.9329,-2.0830) (-0.8830,-2.1696) (-0.8294,-2.2540) (-0.7725,-2.3362) (-0.7125,-2.4162) (-0.6496,-2.4939) (-0.5840,-2.5695) (-0.5161,-2.6428) (-0.4458,-2.7140) (-0.3735,-2.7830) (-0.2992,-2.8500) (-0.2232,-2.9149) (-0.1455,-2.9779)};
\draw[main!65,line width=.35pt] plot[smooth] coordinates {(-0.1200,-1.0000) (-0.1449,-1.0021) (-0.2433,-1.0193) (-0.3373,-1.0530) (-0.4227,-1.1048) (-0.4948,-1.1737) (-0.5506,-1.2565) (-0.5890,-1.3487) (-0.6108,-1.4461) (-0.6175,-1.5458) (-0.6112,-1.6456) (-0.5938,-1.7440) (-0.5668,-1.8403) (-0.5317,-1.9339) (-0.4895,-2.0245) (-0.4413,-2.1121) (-0.3878,-2.1966) (-0.3297,-2.2780) (-0.2676,-2.3563) (-0.2020,-2.4318) (-0.1333,-2.5044) (-0.0617,-2.5743) (0.0122,-2.6415) (0.0884,-2.7063) (0.1665,-2.7688) (0.2464,-2.8289) (0.3279,-2.8869) (0.4107,-2.9429) (0.4949,-2.9969)};
\draw[main!65,line width=.35pt] plot[smooth] coordinates {(-0.1039,-1.0600) (-0.1236,-1.0754) (-0.1933,-1.1469) (-0.2462,-1.2315) (-0.2808,-1.3252) (-0.2973,-1.4237) (-0.2975,-1.5236) (-0.2835,-1.6225) (-0.2577,-1.7191) (-0.2218,-1.8124) (-0.1776,-1.9020) (-0.1265,-1.9880) (-0.0695,-2.0701) (-0.0076,-2.1486) (0.0586,-2.2236) (0.1283,-2.2953) (0.2011,-2.3638) (0.2767,-2.4293) (0.3545,-2.4921) (0.4344,-2.5522) (0.5161,-2.6098) (0.5994,-2.6652) (0.6841,-2.7184) (0.7700,-2.7695) (0.8570,-2.8187) (0.9451,-2.8661) (1.0340,-2.9118) (1.1238,-2.9559) (1.2143,-2.9985)};
\draw[main!65,line width=.35pt] plot[smooth] coordinates {(-0.0600,-1.1039) (-0.0689,-1.1273) (-0.0923,-1.2243) (-0.0968,-1.3241) (-0.0848,-1.4233) (-0.0589,-1.5198) (-0.0216,-1.6125) (0.0250,-1.7009) (0.0792,-1.7850) (0.1395,-1.8647) (0.2050,-1.9403) (0.2746,-2.0120) (0.3479,-2.0801) (0.4241,-2.1448) (0.5028,-2.2065) (0.5837,-2.2652) (0.6665,-2.3213) (0.7509,-2.3749) (0.8368,-2.4262) (0.9239,-2.4753) (1.0120,-2.5225) (1.1012,-2.5678) (1.1912,-2.6113) (1.2820,-2.6532) (1.3735,-2.6935) (1.4657,-2.7324) (1.5584,-2.7699) (1.6516,-2.8062) (1.7452,-2.8412) (1.8393,-2.8750) (1.9338,-2.9077) (2.0287,-2.9395) (2.1238,-2.9702) (2.2193,-2.9999)};
\draw[main!65,line width=.35pt] plot[smooth] coordinates {(-0.0000,-1.1200) (0.0034,-1.1448) (0.0267,-1.2419) (0.0635,-1.3348) (0.1111,-1.4227) (0.1670,-1.5056) (0.2296,-1.5835) (0.2976,-1.6569) (0.3698,-1.7260) (0.4456,-1.7912) (0.5243,-1.8529) (0.6054,-1.9114) (0.6886,-1.9669) (0.7735,-2.0197) (0.8599,-2.0700) (0.9476,-2.1181) (1.0364,-2.1640) (1.1262,-2.2080) (1.2169,-2.2502) (1.3083,-2.2906) (1.4005,-2.3295) (1.4932,-2.3669) (1.5865,-2.4030) (1.6803,-2.4377) (1.7745,-2.4712) (1.8691,-2.5035) (1.9641,-2.5347) (2.0595,-2.5649) (2.1551,-2.5941) (2.2510,-2.6223) (2.3472,-2.6497) (2.4436,-2.6762) (2.5403,-2.7019) (2.6371,-2.7268) (2.7341,-2.7510) (2.8314,-2.7745) (2.9287,-2.7973) (3.0262,-2.8195) (3.1239,-2.8410) (3.2217,-2.8620) (3.3196,-2.8824) (3.4176,-2.9022) (3.5157,-2.9215) (3.6139,-2.9403) (3.7122,-2.9586) (3.8106,-2.9764) (3.9091,-2.9938)};
\draw[main!65,line width=.35pt] plot[smooth] coordinates {(0.0600,-1.1039) (0.0742,-1.1245) (0.1359,-1.2031) (0.2039,-1.2765) (0.2766,-1.3450) (0.3532,-1.4093) (0.4329,-1.4698) (0.5150,-1.5269) (0.5991,-1.5808) (0.6850,-1.6320) (0.7724,-1.6807) (0.8610,-1.7271) (0.9506,-1.7714) (1.0412,-1.8137) (1.1327,-1.8542) (1.2248,-1.8931) (1.3176,-1.9304) (1.4110,-1.9662) (1.5048,-2.0006) (1.5992,-2.0338) (1.6939,-2.0658) (1.7890,-2.0966) (1.8845,-2.1264) (1.9803,-2.1551) (2.0764,-2.1828) (2.1727,-2.2097) (2.2693,-2.2356) (2.3661,-2.2607) (2.4631,-2.2851) (2.5602,-2.3087) (2.6576,-2.3315) (2.7551,-2.3537) (2.8528,-2.3752) (2.9506,-2.3960) (3.0485,-2.4163) (3.1465,-2.4359) (3.2447,-2.4551) (3.3430,-2.4736) (3.4413,-2.4917) (3.5398,-2.5093) (3.6383,-2.5264) (3.7369,-2.5430) (3.8356,-2.5592) (3.9343,-2.5750)};
\draw[main!65,line width=.35pt] plot[smooth] coordinates {(0.1039,-1.0600) (0.1257,-1.0722) (0.2134,-1.1203) (0.3019,-1.1668) (0.3912,-1.2119) (0.4812,-1.2554) (0.5719,-1.2975) (0.6633,-1.3382) (0.7553,-1.3774) (0.8478,-1.4153) (0.9409,-1.4519) (1.0344,-1.4872) (1.1285,-1.5213) (1.2229,-1.5542) (1.3177,-1.5860) (1.4129,-1.6167) (1.5084,-1.6464) (1.6042,-1.6751) (1.7002,-1.7028) (1.7966,-1.7296) (1.8932,-1.7555) (1.9900,-1.7806) (2.0870,-1.8048) (2.1842,-1.8283) (2.2816,-1.8511) (2.3791,-1.8731) (2.4768,-1.8944) (2.5746,-1.9151) (2.6726,-1.9352) (2.7707,-1.9547) (2.8689,-1.9735) (2.9672,-1.9919) (3.0656,-2.0096) (3.1641,-2.0269) (3.2627,-2.0437) (3.3613,-2.0600) (3.4601,-2.0759) (3.5589,-2.0913) (3.6577,-2.1063) (3.7567,-2.1208) (3.8557,-2.1350) (3.9547,-2.1488)};
\draw[main!65,line width=.35pt] plot[smooth] coordinates {(-3.9500,-2.8000) (-3.9263,-2.8080) (-3.8318,-2.8408) (-3.7377,-2.8746) (-3.6441,-2.9096) (-3.5508,-2.9457) (-3.4580,-2.9828)};
\draw[main!65,line width=.35pt] plot[smooth] coordinates {(-3.9500,-2.4000) (-3.9262,-2.4078) (-3.8315,-2.4397) (-3.7371,-2.4729) (-3.6433,-2.5074) (-3.5499,-2.5432) (-3.4571,-2.5804) (-3.3648,-2.6189) (-3.2731,-2.6587) (-3.1819,-2.6998) (-3.0914,-2.7423) (-3.0014,-2.7859) (-2.9120,-2.8307) (-2.8232,-2.8767) (-2.7350,-2.9237) (-2.6472,-2.9718)};
\draw[main!65,line width=.35pt] plot[smooth] coordinates {(-3.9500,-2.0000) (-3.9261,-2.0073) (-3.8307,-2.0373) (-3.7358,-2.0689) (-3.6414,-2.1019) (-3.5476,-2.1365) (-3.4544,-2.1727) (-3.3618,-2.2105) (-3.2699,-2.2499) (-3.1788,-2.2910) (-3.0883,-2.3338) (-2.9987,-2.3781) (-2.9099,-2.4240) (-2.8218,-2.4714) (-2.7346,-2.5202) (-2.6481,-2.5704) (-2.5623,-2.6219) (-2.4773,-2.6745) (-2.3929,-2.7281) (-2.3091,-2.7827) (-2.2258,-2.8380) (-2.1429,-2.8940) (-2.0604,-2.9506) (-1.9782,-3.0075)};
\draw[main!65,line width=.35pt] plot[smooth] coordinates {(-3.9500,-1.6000) (-3.9259,-1.6065) (-3.8296,-1.6334) (-3.7337,-1.6619) (-3.6384,-1.6921) (-3.5436,-1.7240) (-3.4495,-1.7578) (-3.3560,-1.7934) (-3.2634,-1.8311) (-3.1716,-1.8707) (-3.0807,-1.9124) (-2.9908,-1.9561) (-2.9019,-2.0019) (-2.8140,-2.0497) (-2.7273,-2.0995) (-2.6417,-2.1512) (-2.5572,-2.2047) (-2.4738,-2.2598) (-2.3914,-2.3164) (-2.3099,-2.3744) (-2.2293,-2.4335) (-2.1494,-2.4937) (-2.0701,-2.5547) (-1.9914,-2.6164) (-1.9131,-2.6785) (-1.8350,-2.7410) (-1.7571,-2.8037) (-1.6791,-2.8663) (-1.6011,-2.9289) (-1.5230,-2.9913)};
\draw[main!65,line width=.35pt] plot[smooth] coordinates {(-3.9500,-1.2000) (-3.9256,-1.2054) (-3.8281,-1.2277) (-3.7310,-1.2515) (-3.6343,-1.2770) (-3.5381,-1.3043) (-3.4425,-1.3335) (-3.3475,-1.3648) (-3.2533,-1.3983) (-3.1599,-1.4342) (-3.0675,-1.4724) (-2.9763,-1.5133) (-2.8862,-1.5567) (-2.7975,-1.6029) (-2.7102,-1.6517) (-2.6245,-1.7032) (-2.5404,-1.7573) (-2.4579,-1.8138) (-2.3770,-1.8726) (-2.2977,-1.9335) (-2.2198,-1.9962) (-2.1433,-2.0606) (-2.0680,-2.1264) (-1.9936,-2.1933) (-1.9201,-2.2610) (-1.8471,-2.3293) (-1.7745,-2.3981) (-1.7020,-2.4671) (-1.6296,-2.5360) (-1.5570,-2.6048) (-1.4842,-2.6733) (-1.4108,-2.7413) (-1.3370,-2.8087) (-1.2625,-2.8754) (-1.1873,-2.9413) (-1.1114,-3.0064)};
\draw[main!65,line width=.35pt] plot[smooth] coordinates {(-3.9500,-0.8000) (-3.9253,-0.8038) (-3.8266,-0.8199) (-3.7281,-0.8373) (-3.6299,-0.8560) (-3.5320,-0.8763) (-3.4344,-0.8984) (-3.3374,-0.9223) (-3.2408,-0.9484) (-3.1449,-0.9768) (-3.0498,-1.0077) (-2.9557,-1.0415) (-2.8627,-1.0783) (-2.7711,-1.1184) (-2.6811,-1.1619) (-2.5930,-1.2091) (-2.5069,-1.2600) (-2.4231,-1.3146) (-2.3419,-1.3729) (-2.2632,-1.4346) (-2.1871,-1.4995) (-2.1135,-1.5672) (-2.0423,-1.6374) (-1.9732,-1.7096) (-1.9058,-1.7835) (-1.8397,-1.8586) (-1.7747,-1.9346) (-1.7104,-2.0111) (-1.6463,-2.0879) (-1.5823,-2.1647) (-1.5179,-2.2413) (-1.4531,-2.3174) (-1.3876,-2.3929) (-1.3211,-2.4677) (-1.2537,-2.5416) (-1.1853,-2.6144) (-1.1156,-2.6862) (-1.0448,-2.7568) (-0.9728,-2.8262) (-0.8995,-2.8943) (-0.8251,-2.9610)};
\draw[main!65,line width=.35pt] plot[smooth] coordinates {(-3.9500,-0.4000) (-3.9251,-0.4020) (-3.8254,-0.4105) (-3.7259,-0.4197) (-3.6264,-0.4297) (-3.5270,-0.4406) (-3.4277,-0.4526) (-3.3286,-0.4658) (-3.2296,-0.4804) (-3.1310,-0.4966) (-3.0326,-0.5146) (-2.9347,-0.5348) (-2.8372,-0.5574) (-2.7405,-0.5828) (-2.6448,-0.6116) (-2.5502,-0.6441) (-2.4573,-0.6810) (-2.3665,-0.7228) (-2.2783,-0.7700) (-2.1936,-0.8231) (-2.1129,-0.8821) (-2.0368,-0.9469) (-1.9657,-1.0172) (-1.8997,-1.0923) (-1.8383,-1.1713) (-1.7811,-1.2533) (-1.7273,-1.3375) (-1.6759,-1.4233) (-1.6261,-1.5100) (-1.5771,-1.5972) (-1.5283,-1.6845) (-1.4790,-1.7715) (-1.4290,-1.8581) (-1.3777,-1.9439) (-1.3249,-2.0289) (-1.2704,-2.1127) (-1.2142,-2.1954) (-1.1560,-2.2767) (-1.0959,-2.3567) (-1.0339,-2.4351) (-0.9699,-2.5119) (-0.9040,-2.5871) (-0.8362,-2.6607) (-0.7667,-2.7325) (-0.6954,-2.8027) (-0.6225,-2.8711) (-0.5480,-2.9378) (-0.4721,-3.0029)};
\draw[main!65,line width=.35pt] plot[smooth] coordinates {(-3.9500,0.0000) (-3.9250,0.0000) (-3.8250,0.0000) (-3.7250,0.0000) (-3.6250,0.0000) (-3.5250,0.0000) (-3.4250,0.0000) (-3.3250,0.0000) (-3.2250,0.0000) (-3.1250,0.0000) (-3.0250,0.0000) (-2.9250,0.0000) (-2.8250,0.0000) (-2.7250,0.0000) (-2.6250,0.0000) (-2.5250,0.0000) (-2.4250,0.0000) (-2.3250,0.0000) (-2.2250,0.0000) (-2.1250,0.0000) (-2.0250,0.0000) (-1.9250,0.0000) (-1.8250,0.0000) (-1.7250,0.0000) (-1.6250,0.0000) (-1.5250,0.0000) (-1.4250,0.0000) (-1.3250,0.0000) (-1.2250,0.0000) (-1.1250,0.0000) (-1.0250,0.0000)};
\draw[main!65,line width=.35pt] plot[smooth] coordinates {(-3.9500,0.4000) (-3.9251,0.4020) (-3.8254,0.4105) (-3.7259,0.4197) (-3.6264,0.4297) (-3.5270,0.4406) (-3.4277,0.4526) (-3.3286,0.4658) (-3.2296,0.4804) (-3.1310,0.4966) (-3.0326,0.5146) (-2.9347,0.5348) (-2.8372,0.5574) (-2.7405,0.5828) (-2.6448,0.6116) (-2.5502,0.6441) (-2.4573,0.6810) (-2.3665,0.7228) (-2.2783,0.7700) (-2.1936,0.8231) (-2.1129,0.8821) (-2.0368,0.9469) (-1.9657,1.0172) (-1.8997,1.0923) (-1.8383,1.1713) (-1.7811,1.2533) (-1.7273,1.3375) (-1.6759,1.4233) (-1.6261,1.5100) (-1.5771,1.5972) (-1.5283,1.6845) (-1.4790,1.7715) (-1.4290,1.8581) (-1.3777,1.9439) (-1.3249,2.0289) (-1.2704,2.1127) (-1.2142,2.1954) (-1.1560,2.2767) (-1.0959,2.3567) (-1.0339,2.4351) (-0.9699,2.5119) (-0.9040,2.5871) (-0.8362,2.6607) (-0.7667,2.7325) (-0.6954,2.8027) (-0.6225,2.8711) (-0.5480,2.9378) (-0.4721,3.0029)};
\draw[main!65,line width=.35pt] plot[smooth] coordinates {(-3.9500,0.8000) (-3.9253,0.8038) (-3.8266,0.8199) (-3.7281,0.8373) (-3.6299,0.8560) (-3.5320,0.8763) (-3.4344,0.8984) (-3.3374,0.9223) (-3.2408,0.9484) (-3.1449,0.9768) (-3.0498,1.0077) (-2.9557,1.0415) (-2.8627,1.0783) (-2.7711,1.1184) (-2.6811,1.1619) (-2.5930,1.2091) (-2.5069,1.2600) (-2.4231,1.3146) (-2.3419,1.3729) (-2.2632,1.4346) (-2.1871,1.4995) (-2.1135,1.5672) (-2.0423,1.6374) (-1.9732,1.7096) (-1.9058,1.7835) (-1.8397,1.8586) (-1.7747,1.9346) (-1.7104,2.0111) (-1.6463,2.0879) (-1.5823,2.1647) (-1.5179,2.2413) (-1.4531,2.3174) (-1.3876,2.3929) (-1.3211,2.4677) (-1.2537,2.5416) (-1.1853,2.6144) (-1.1156,2.6862) (-1.0448,2.7568) (-0.9728,2.8262) (-0.8995,2.8943) (-0.8251,2.9610)};
\draw[main!65,line width=.35pt] plot[smooth] coordinates {(-3.9500,1.2000) (-3.9256,1.2054) (-3.8281,1.2277) (-3.7310,1.2515) (-3.6343,1.2770) (-3.5381,1.3043) (-3.4425,1.3335) (-3.3475,1.3648) (-3.2533,1.3983) (-3.1599,1.4342) (-3.0675,1.4724) (-2.9763,1.5133) (-2.8862,1.5567) (-2.7975,1.6029) (-2.7102,1.6517) (-2.6245,1.7032) (-2.5404,1.7573) (-2.4579,1.8138) (-2.3770,1.8726) (-2.2977,1.9335) (-2.2198,1.9962) (-2.1433,2.0606) (-2.0680,2.1264) (-1.9936,2.1933) (-1.9201,2.2610) (-1.8471,2.3293) (-1.7745,2.3981) (-1.7020,2.4671) (-1.6296,2.5360) (-1.5570,2.6048) (-1.4842,2.6733) (-1.4108,2.7413) (-1.3370,2.8087) (-1.2625,2.8754) (-1.1873,2.9413) (-1.1114,3.0064)};
\draw[main!65,line width=.35pt] plot[smooth] coordinates {(-3.9500,1.6000) (-3.9259,1.6065) (-3.8296,1.6334) (-3.7337,1.6619) (-3.6384,1.6921) (-3.5436,1.7240) (-3.4495,1.7578) (-3.3560,1.7934) (-3.2634,1.8311) (-3.1716,1.8707) (-3.0807,1.9124) (-2.9908,1.9561) (-2.9019,2.0019) (-2.8140,2.0497) (-2.7273,2.0995) (-2.6417,2.1512) (-2.5572,2.2047) (-2.4738,2.2598) (-2.3914,2.3164) (-2.3099,2.3744) (-2.2293,2.4335) (-2.1494,2.4937) (-2.0701,2.5547) (-1.9914,2.6164) (-1.9131,2.6785) (-1.8350,2.7410) (-1.7571,2.8037) (-1.6791,2.8663) (-1.6011,2.9289) (-1.5230,2.9913)};
\draw[main!65,line width=.35pt] plot[smooth] coordinates {(-3.9500,2.0000) (-3.9261,2.0073) (-3.8307,2.0373) (-3.7358,2.0689) (-3.6414,2.1019) (-3.5476,2.1365) (-3.4544,2.1727) (-3.3618,2.2105) (-3.2699,2.2499) (-3.1788,2.2910) (-3.0883,2.3338) (-2.9987,2.3781) (-2.9099,2.4240) (-2.8218,2.4714) (-2.7346,2.5202) (-2.6481,2.5704) (-2.5623,2.6219) (-2.4773,2.6745) (-2.3929,2.7281) (-2.3091,2.7827) (-2.2258,2.8380) (-2.1429,2.8940) (-2.0604,2.9506) (-1.9782,3.0075)};
\draw[main!65,line width=.35pt] plot[smooth] coordinates {(-3.9500,2.4000) (-3.9262,2.4078) (-3.8315,2.4397) (-3.7371,2.4729) (-3.6433,2.5074) (-3.5499,2.5432) (-3.4571,2.5804) (-3.3648,2.6189) (-3.2731,2.6587) (-3.1819,2.6998) (-3.0914,2.7423) (-3.0014,2.7859) (-2.9120,2.8307) (-2.8232,2.8767) (-2.7350,2.9237) (-2.6472,2.9718)};
\draw[main!65,line width=.35pt] plot[smooth] coordinates {(-3.9500,2.8000) (-3.9263,2.8080) (-3.8318,2.8408) (-3.7377,2.8746) (-3.6441,2.9096) (-3.5508,2.9457) (-3.4580,2.9828)};
\draw[->,gray](-4.1,0)--(4.45,0) node[right]{$x$};
\draw[->,gray](0,-3.1)--(0,3.15) node[above]{$y$};
\fill (0,1)circle(2pt) node[above right,fill=white,inner sep=1pt]{源 $\ii$};
\fill (0,-1)circle(2pt) node[below right,fill=white,inner sep=1pt]{源 $-\ii$};
\fill (-1,0)circle(2pt) node[below left,fill=white,inner sep=1pt]{$-1$};
\draw[->,thick](3.0000,2.6000)--(3.5375,2.7164);
\draw[->,thick](3.0000,-2.6000)--(3.5375,-2.7164);
\draw[->,thick](0.2200,1.0000)--(0.5188,1.0264);
\draw[->,thick](0.2200,-1.0000)--(0.5188,-1.0264);
\end{tikzpicture}\qedhere\end{center}

\end{solution}
\begin{exercise}\label{ass:pex2-12}
\textbf{原题 Pex2.12。} 求 (a) $\int_{-\pi}^{\pi}(1+\sin^2\theta)^{-1}\dd\theta$（原题附答案 $\pi\sqrt2$）; (b) $\int_{-\infty}^{\infty}x/(x^2+4x+13)^2\dd x$（原题附答案 $-\pi/27$）; (c) $\mathrm{p.v.}\int_{-\infty}^{\infty}x\sin x/(1+x^2)\dd x$; (d) $\mathrm{p.v.}\int_{-\infty}^{\infty}\cos x/(x+\ii)\dd x$; (e) $\mathrm{p.v.}\int_{-\infty}^{\infty}x\ee^{2\ii x}/(x^2-1)\dd x$.
\end{exercise}
\begin{solution}
据官方解答翻译并补全弧消失依据. (a) 取 $z=\ee^{\ii\theta}$, 积分变为 $\int_{|z|=1}-4z/(\ii(z^4-6z^2+1))\dd z$. 圆内简单极点为 $\pm(\sqrt2-1)$, 每点留数 $1/(\ii\sqrt8)$, 故值 $2\pi\ii\cdot2/(\ii\sqrt8)=\pi\sqrt2$.

(b) 上半平面只有二阶极点 $a=-2+3\ii$, 下方另一个为 $b=-2-3\ii$. 其留数为 $[z/(z-b)^2]'|_a=(-a-b)/(a-b)^3=\ii/54$. 大弧上为 $O(R^{-3})$, 弧积分消失, 得 $2\pi\ii\cdot\ii/54=-\pi/27$.

(c) 改为 $\widetilde I=\mathrm{p.v.}\int x\ee^{\ii x}/(1+x^2)\dd x$, 取虚部. 上半平面极点 $\ii$ 留数 $\ee^{-1}/2$, Jordan 弧估计适用, 因有理因子为 $O(1/R)$. 故 $\widetilde I=\pi\ii/\ee$, 所求 $\pi/\ee$.

(d) 把余弦分为 $(\ee^{\ii x}+\ee^{-\ii x})/2$. 第一项在上半平面闭合, 无极点得0; 第二项在下半平面顺时针闭合, $-\ii$ 处留数 $\ee^{-1}$, 得 $-2\pi\ii/\ee$. 各弧由 Jordan 估计消失; 除以2得 $-\pi\ii/\ee$.

(e) 同时在 $\pm1$ 左右对称挖去小邻域定义主值, 并在上方用顺时针小半圆避开极点. 大弧消失, 域内无极点. 每个小弧贡献趋于 $-\pi\ii$ 乘相应留数; 留数为 $\ee^{-2\ii}/2$、$\ee^{2\ii}/2$. 因而主值为 $\pi\ii(\ee^{-2\ii}+\ee^{2\ii})/2=\pi\ii\cos2$. 这些主值符号不能被理解为在实极点上直接作通常反常积分.
\end{solution}
